\section{Introduction}

B.17 established that parameter gaming is bounded ($< 0.3$ seats nationally);
R.2 established that input data manipulation is detectable via SHA-256.
R.3 examines the third and most subtle gaming vector: geographic definition choices
that appear to be neutral technical decisions but may be strategically chosen.

Two geographic decisions are relevant:
\begin{enumerate}
  \item \textbf{Resolution choice}: whether to use census tracts or block groups
    as the atomic unit. F.3 establishes a $k/n > 0.05$ threshold rule, but
    an adversary might argue for a different threshold to shift a state from
    one resolution level to another.
  \item \textbf{Adjacency definition}: whether two geographically adjacent tracts
    share an adjacency edge. The standard definition (any shared TIGER/Line boundary)
    is objective, but an adversary might use a different adjacency criterion
    (e.g., require minimum shared boundary length) to exclude specific cross-partisan
    edges.
\end{enumerate}

This paper tests the maximum partisan shift achievable via each vector and shows
that the audit chain makes the manipulation detectable.
